\section{¿Qué es Denait Hartenberg?}
\textbf{Denavid Hartenberg} (D-H) es un método y convención en robótica que asigna \textbf{marcos de referencia} a \textbf{eslabones} de una \textbf{cadena cinemática abierta simple}. Además, permite hacer \textbf{transformaciones} entre \textbf{marcos de referencia}, es decir, permite ubicar a un \textbf{marco de referencia} respecto de otro. En general, D-H nos permite modelar una cadena cinemática abierta simple, como un brazo robótico.

Existen otras convenciones, hay muchas formas y orientaciones para acomodar un marco de referencia (e.g. utilizar el sistema de la mano izquierda en lugar del de la mano derecha), hay diferentes ordenes para rotar y trasladar marcos de referencia. Al profundizar en diferentes recursos externos, es importante saber que pueden utilizar una convencion diferente para no confundirnos.